\problemname{Bisecting Bargain}

\illustration{0.3}{illustration-small.jpg}{
    Illustration of Sample Output 4. \\
    Photo by Jeroen Op de Beek
}

Emilia and Alex love Christmas markets.
There are so many stalls to explore and so many delicious foods to try!
The options seem endless: Schupfnudeln with Sauerkraut, langos, crêpes, roasted almonds, Bratwurst, twister fries and much more.

Here in Germany, some stalls still do not accept credit cards, so Emilia and
Alex need to withdraw some cash from a nearby cash machine (also called an ATM,
or Automated Teller Machine).
Due to some fees, it is cheaper to withdraw all the cash they need at once, so they plan to do so.

This particular cash machine works in the following way: the user inputs an integer $x$, and then the machine chooses a selection of coins and banknotes that sums to €$x$.
The cash machine can dispense all possible Euro coins and notes with values €$1$ and up: €$1$, €$2$, €$5$, €$10$, €$20$, €$50$, €$100$, €$200$, and €$500$.
It has sufficiently many coins and notes of each type to form any combination that
sums to $x$, and
you do not know in advance which of these combinations it will dispense.


Since Emilia and Alex plan to visit some different stalls independently, they want to split the withdrawn amount evenly between them.
While waiting in the long queue in front of the cash machine, Emilia suddenly realizes that this might not be possible, depending on which coins and banknotes the cash machine dispenses.
For example, €$42$ might not be evenly splittable (see Sample 1),
whereas no matter how €$40$ is dispensed by the cash machine, the cash can
always be divided into two piles of €$20$.



% Depending on what they encounter at the Christmas market, Alex and Emily want to withdraw various amounts.
% Help them check whether $x$€ dispensed by the cash machine can always be split evenly.
Can the money always be split evenly if they withdraw €$n$ from the cash machine? 

% Alex and Emily are at a Christmas market in front of an cash machine and want to withdraw $42$€, when suddenly Alex asks:
% ``If we do this, are you sure that we can evenly split the money between the two of us?''
% Emily thinks for a short moment and realizes that this indeed could be impossible if the cash machine dispenses some specific coins.
% However, she is also sure that if she only withdraws $40$€, they can always split it evenly, regardless of what the cash machine would dispense.

% The cash machine can dispense all possible Euro coins and notes, $1$€ and up, i.e., $1$€, $2$€, $5$€, $10$€, $20$€, $50$€, $100$€, $200$€, and $500$€.
% Depending on what they encounter at the Christmas market, Alex and Emily want to withdraw various amounts.
% Help them check whether $x$€ dispensed by the cash machine can always be split evenly.



\begin{Input}
    The input consists of:
    \begin{itemize}
        \item One line with an integer $n$ ($1\leq n\leq 10\,000$), the amount of
    cash in € that Emilia and Alex want to withdraw in total.
    \end{itemize}
\end{Input}

\begin{Output}
    If the money can always be split evenly, output ``\texttt{splittable}''.
    Otherwise, output the number of
    coins and notes, followed by their values, such that the values add to $n$
    and the money cannot be split evenly.

    If there are multiple ways of choosing coins and notes that are not splittable, you may output any one of them.
\end{Output}
